% Part: proof-theory
% Chapter: natural-deduction
% Section: sequents

\documentclass[../../../include/open-logic-section]{subfiles}

\begin{document}

\olfileid{pt}{nat}{seq}

\olsection{Deduksi Alami nganggo Sekuèn: \Log{N2c} lan~\Log{N2i}}

!!^a{proof}~$\delta$ ing \Log{N1c} utawa \Log{N1i} nuduhaké yèn
!!{formula} pungkasané~$!A$ ndhèrèk saka asumsi kabukané. Yèn
$\Gamma$ iku himpunan !!{formula}s~$!B$ kang $!B^x$ katon minangka
asumsi kabuka ing~$\delta$, kita bisa nulis iki minangka
$\delta \Proves \Gamma \Sequent !A$. Iki mènèhi gagasan yèn
sistem deduksi alami uga bisa dirumusaké kanggo sekuèn. Ing sistem
kaya mangkono, saben sekuèn nggawa katrangan bab asumsi kabuka endi
kang dadi dhasaré !!{formula} ing sisih tengen, yaiku asumsi kang
kabuka ing sub-!!{proof} kang pungkasané sekuèn mau. Iki deduksi
alami nganggo sekuèn, utawa deduksi alami ``gaya sekuèn,'' lan
sistem asilé diarani \Log{N2c} lan~\Log{N2i}.

Supaya korespondènsi antarané \Log{N1} lan \Log{N2} luwih cedhak,
kita gawe himpunan~$\Gamma$ ing sisih kiwa isi !!{formula}s
mawa label~$x:!B$. Amarga aturan ing \Log{N1} mung nggarap
!!{formula}s ing premis, yaiku !!{formula}s pungkasan saka
sub-!!{proof}s langsung, aturan \Log{N2} mung nggarap
konsekuèn sekuèn. Mula formula mawa label ing sisih kiwa
sekuèn bisa diarani kontèks.

Ing panggonané asumsi, bukti ing \Log{N2} nduwèni
aksioma minangka sekuèn paling dhuwur kanthi wujud
\[x:!A \Sequent !A.\]
Aturané cocog persis karo aturan \Log{N1} lan katulis ing
\olref{tab:N2}.

\subfile{rules-N2}

Amarga kita kudu ngidinaké (kaya ing \Log{N1c} lan \Log{N1i})
yèn aturan \Intro{\lif}, \Intro{\lnot}, \Elim{\lor}, lan \Elim{\lexists}
bisa mbusak asumsi nanging ora kudu nglakoni, panrapan aturan kang
cocog ing \Log{N2c} lan~\Log{N2i} uga dianggep bener sanajan
!!{formula}s kontèks kang cocog, $!A$, $!B$, lan $!A(a)$,
ora ana ing kontèks premis kang gegayutan.

\begin{prop}
Yèn $\delta$ iku !!a{proof} ing~\Log{N1c} utawa \Log{N1i}
kanggo~$!A$, banjur ana !!a{proof}~$\delta'$ kanggo
$\Gamma \Sequent !A$ ing~\Log{N2c} (\Log{N2i}), kanthi $\Gamma$
minangka himpunan kabèh~$x:!B$ saéngga $!B^x$ iku asumsi kabuka
saka~$\delta$.
\end{prop}

\begin{proof}
  Kanthi induksi marang $n=\pheight{\delta}$. Yèn $n=0$,
  $\delta$ mung isi siji asumsi kabuka~$!A^x$. Mula
  $\Gamma = \{x:!A\}$ lan $\Gamma \Sequent !A$ iku aksioma
  ing \Log{N2c} utawa~\Log{N2i}.

  Yèn $n>0$, kita mbedakaké kasus miturut inferènsi pungkasan
  ing~$\delta$. Kita mung ngrembug siji kasus; sisané dadi
  gladhen.

  Upamana inferènsi pungkasané $\Intro{\lif}$. Banjur $\delta$
  pungkasané kaya mangkéné:
  \[
  \AxiomC{$\Discharge{!B}{x}$}
  \RightLabel{$\delta_1$}
  \DeduceC{$!C$}
  \DischargeRule{\Intro{\lif}}{x}
  \UnaryInfC{$!B \lif !C$}
  \DisplayProof
  \]
  Miturut hipotesis induksi ana !!a{proof}~$\delta_1'$ ing~\Log{N2c}
  (\Log{N2i}) kanggo $\Gamma' \Sequent !C$, kanthi $\Gamma'$
  minangka asumsi mawa label kang kabuka ing~$\delta_1$.
  Yèn $!B^x$ iku asumsi kabuka saka $\delta_1$, asumsi iku dibusak
  ing inferènsi \Intro{\lif}, mula asumsi kabuka saka~$\delta$
  yaiku $\Gamma = \Gamma' \setminus \{x:!B\}$. Tegesé,
  $\delta_1'$ iku !!a{proof} kanggo~$x:!B, \Gamma \Sequent !C$,
  lan $\delta'$ bisa dipilih dadi
  \[
  \AxiomC{}
  \RightLabel{$\delta_1'$}
  \Deduce$x:!B, \Gamma \fCenter !C$
  \RightLabel{\Intro{\lif}}
  \UnaryInf$\Gamma \fCenter !B \lif !C$
  \DisplayProof
  \]
  Yèn $!B^x$ dudu asumsi kabuka saka~$\delta_1$, banjur
  \Intro{\lif} ora mbusak asumsi, lan asumsi kabuka saka $\delta$
  lan $\delta_1$ padha. Amarga formula kontèks $x:!B$ ora kudu
  ana ing premis \Intro{\lif} ing~\Log{N2c} lan~\Log{N2i},
  $\delta'$ bisa dipilih dadi
  \[
  \AxiomC{}
  \RightLabel{$\delta_1'$}
  \Deduce$\Gamma \fCenter !C$
  \RightLabel{\Intro{\lif}}
  \UnaryInf$\Gamma \fCenter !B \lif !C$
  \DisplayProof
  \]

  Kasus nalika $\delta$ pungkasané aturan liya ditangani kanthi
  cara kang padha.
\end{proof}

\begin{prop}
Yèn $\delta$ iku !!a{proof} ing \Log{N2c} utawa \Log{N2i}
kanggo $\Gamma \Sequent !A$, banjur ana !!a{proof}~$\delta'$
ing \Log{N1c} (\Log{N1i}) kanthi !!{formula} pungkasan~$!A$
lan asumsi kabuka $!B^x$ kanggo saben $x:!B \in \Gamma$.
\end{prop}

\begin{proof}
  Kita nganggo induksi manèh marang !!{height}~$n$ saka~$\delta$.
  Yèn $n=0$, $\delta$ iku aksioma, $x:!A \Sequent !A$.
  !!{proof}~$\delta'$ yaiku asumsi~$!A^x$.

  Yèn $n>0$, kita mbedakaké kasus miturut inferènsi pungkasan
  ing~$\delta$. Tuladhané, yèn inferènsi iku~$\Intro{\lif}$,
  $\delta$~pungkasané kaya mangkéné:
  \[
  \bottomAlignProof
  \AxiomC{}
  \RightLabel{$\delta_1$}
  \Deduce$x:!B, \Gamma \fCenter !C$
  \RightLabel{\Intro{\lif}}
  \UnaryInf$\Gamma \fCenter !B \lif !C$
  \DisplayProof
  \quad\text{ utawa }\quad
  \bottomAlignProof
  \AxiomC{}
  \RightLabel{$\delta_1$}
  \Deduce$\Gamma \fCenter !C$
  \RightLabel{\Intro{\lif}}
  \UnaryInf$\Gamma \fCenter !B \lif !C$
  \DisplayProof
  \]
  Miturut hipotesis induksi, ana !!a{proof} ing~\Log{N1c}
  (\Log{N1i}) kanggo~$!C$. Ing kasus kapisan, $!B^x$ iku
  asumsi kabuka ing subbukti ndhuwur~$\delta'$;
  ing kasus kapindho ora. !!{proof}~$\delta'$ bisa dipilih dadi
  \[
  \bottomAlignProof
  \AxiomC{$\Discharge{!B}{x}$}
  \RightLabel{$\delta_1$}
  \DeduceC{$!C$}
  \DischargeRule{\Intro{\lif}}{x}
  \UnaryInfC{$!B \lif !C$}
  \DisplayProof
  \quad\text{ utawa }\quad
  \bottomAlignProof
  \AxiomC{}
  \RightLabel{$\delta_1$}
  \DeduceC{$!C$}
  \RightLabel{\Intro{\lif}}
  \UnaryInfC{$!B \lif !C$}
  \DisplayProof
  \]
  miturut urutané.
\end{proof}

\end{document}
